%% ============================================================
%% Artigo — Agentic DAST
%% Ester Morais Neves — CEFET-MG — 2026
%% Formato: artigo acadêmico, duas colunas, template SBC
%% Compilar com pdfLaTeX (Overleaf: compilador padrão)
%% ============================================================
\documentclass[
    10pt,
    twocolumn,
    a4paper
]{article}

%% ---- Codificação e idioma ------------------------------------
\usepackage[T1]{fontenc}
\usepackage[utf8]{inputenc}
\usepackage{lmodern}
\usepackage[brazil]{babel}

%% ---- Formatação geral ----------------------------------------
\usepackage[
    top=2.5cm, bottom=2.5cm,
    left=1.5cm, right=1.5cm,
    columnsep=0.8cm
]{geometry}
\usepackage{indentfirst}
\usepackage{microtype}
\usepackage{setspace}
\usepackage{titlesec}

\titlespacing*{\section}{0pt}{8pt}{4pt}
\titlespacing*{\subsection}{0pt}{6pt}{3pt}

%% ---- Tabelas -------------------------------------------------
\usepackage{booktabs}
\usepackage{array}
\usepackage{tabularx}
\usepackage{multirow}

%% ---- Matemática ----------------------------------------------
\usepackage{amsmath}
\usepackage{amssymb}

%% ---- Código-fonte --------------------------------------------
\usepackage{listings}
\usepackage{xcolor}

\definecolor{codeblue}{rgb}{0.0, 0.2, 0.6}
\definecolor{codegray}{rgb}{0.5,0.5,0.5}
\definecolor{backcolor}{rgb}{0.97,0.97,0.97}

\lstdefinestyle{artigostyle}{
    backgroundcolor=\color{backcolor},
    commentstyle=\color{codegray},
    keywordstyle=\bfseries\color{codeblue},
    basicstyle=\scriptsize\ttfamily,
    breaklines=true,
    frame=single,
    numbers=left,
    numberstyle=\tiny\color{codegray},
    numbersep=4pt,
    tabsize=2
}
\lstset{style=artigostyle}

%% ---- Figuras -------------------------------------------------
\usepackage{graphicx}
\usepackage{caption}

%% ---- Hyperlinks ----------------------------------------------
\usepackage{hyperref}
\hypersetup{
    colorlinks=true,
    linkcolor=black,
    citecolor=black,
    urlcolor=blue
}

%% ---- Referências (ABNT) --------------------------------------
\usepackage[alf, abnt-emphasize=bf]{abntex2cite}

%% ---- Abstract bilíngue --------------------------------------
\usepackage{abstract}
\renewcommand{\abstractnamefont}{\normalfont\bfseries\small}
\renewcommand{\abstracttextfont}{\normalfont\small}

%% ============================================================
%% Metadados do artigo
%% ============================================================

\title{
    \textbf{Agentic DAST: Um Framework para Detecção Automatizada
    de Vulnerabilidades em Single Page Applications
    via Orquestração de LLMs e Navegação Semântica}
}

\author{
    Ester Morais Neves\textsuperscript{1},
    Mateus Felipe Tymburibá Ferreira\textsuperscript{1,2}\\
    \small\textsuperscript{1}Centro Federal de Educação Tecnológica de
    Minas Gerais (CEFET-MG), Belo Horizonte, MG, Brasil\\
    \small\textsuperscript{2}Universidade Federal de Minas Gerais (UFMG),
    Belo Horizonte, MG, Brasil\\
    \small\texttt{ester.neves@estudante.cefetmg.br},
    \texttt{tymburiba@cefetmg.br}
}

\date{2026}

%% ============================================================
\begin{document}
%% ============================================================

\maketitle
\thispagestyle{empty}

% =============================================================
% Resumo — escrito por último, seguindo a estrutura de 4 passos:
% (1) Contexto/Problema, (2) Solução, (3) Principal Resultado,
% (4) Impacto/Implicação
% =============================================================

\begin{abstract}
Ferramentas de \textit{Dynamic Application Security Testing} (DAST)
tradicionais falham em dois aspectos críticos: são incapazes de navegar
os estados internos de \textit{Single Page Applications} (SPAs) e não
detectam vulnerabilidades que exigem compreensão semântica do domínio,
como Lógica de Negócio, IDOR, XSS armazenado e ataques a JWT.
Este artigo propõe o \textbf{Agentic DAST}, um \textit{framework}
que integra raciocínio semântico de Grandes Modelos de Linguagem (LLMs)
com automação de navegadores via Puppeteer e padrão Screenplay, operando
em um \textit{loop} ReAct autônomo estruturado em quatro módulos:
Percepção, Cognitivo, Executor e Auditoria.
% [A preencher após os experimentos]
A validação empírica é conduzida sobre o OWASP Juice Shop em experimento
comparativo com três cenários (ZAP tradicional, ZAP guiado por E2E e
Agentic DAST autônomo), avaliando \textit{Recall}, \textit{Precision} e
F1-\textit{Score} por categoria de vulnerabilidade.
Os resultados têm potencial de evidenciar que agentes LLM superam
scanners tradicionais na detecção de vulnerabilidades semânticas em SPAs,
abrindo caminho para integração em \textit{pipelines} de DevSecOps.

\bigskip\noindent
\textbf{Palavras-chave:}
DAST, Single Page Application, LLM, Vulnerabilidade de Lógica de Negócio,
XSS, JWT, ReAct, Screenplay Pattern, OWASP Juice Shop.
\end{abstract}

\selectlanguage{english}
\begin{abstract}
Traditional Dynamic Application Security Testing (DAST) tools fail on
two critical fronts: they cannot navigate the internal states of Single
Page Applications (SPAs), and they cannot detect vulnerabilities requiring
semantic domain understanding, such as Business Logic, IDOR, stored XSS,
and JWT attacks. This paper proposes \textbf{Agentic DAST}, a framework
that integrates semantic reasoning from Large Language Models (LLMs) with
browser automation via Puppeteer and the Screenplay Pattern, operating
within an autonomous ReAct loop structured in four modules: Perception,
Cognitive, Executor, and Audit.
% [To be updated after experiments]
Empirical validation is conducted on OWASP Juice Shop through a
three-scenario comparative experiment (traditional ZAP, E2E-guided ZAP,
and autonomous Agentic DAST), evaluating Recall, Precision, and F1-Score
per vulnerability category. Results are expected to demonstrate that LLM
agents outperform traditional scanners in detecting semantic vulnerabilities
in SPAs, paving the way for integration into DevSecOps pipelines.

\bigskip\noindent
\textbf{Keywords:}
DAST, Single Page Application, LLM, Business Logic Vulnerability,
XSS, JWT, ReAct, Screenplay Pattern, OWASP Juice Shop.
\end{abstract}
\selectlanguage{brazil}

% =============================================================
\section{Introdução}
% =============================================================

% --- Parte 1: Contexto ----------------------------------------
A segurança de aplicações web tornou-se um requisito crítico no ciclo de
desenvolvimento de \textit{software}. O custo médio de uma violação de
dados em 2024 atingiu US\$\,4,88 milhões, segundo o relatório da IBM
\cite{ibm2024}, e a OWASP estima que mais de 70\% das vulnerabilidades
exploradas em produção poderiam ter sido detectadas na fase de testes
\cite{owasp2021top10}. O \textit{Dynamic Application Security Testing}
(DAST) é a categoria de ferramentas que testa aplicações em execução,
simulando um atacante externo sem acesso ao código-fonte.

% --- Parte 2: Problema ----------------------------------------
A transição da web para o paradigma das \textit{Single Page Applications}
(SPAs) --- dominadas por \textit{frameworks} como Angular, React e
Vue.js --- criou uma lacuna crítica que as ferramentas DAST consolidadas
não conseguem endereçar \cite{brightsec2024, stackhawk2024}. Em uma SPA,
toda a navegação é controlada por JavaScript no lado do cliente: não há
recarregamento de página, os estados internos (painéis administrativos,
fluxos de \textit{checkout}, carrinhos de compras) são alcançados por
interações com o DOM, e a superfície de ataque real só é revelada após
execução de código JavaScript. Ferramentas como o OWASP ZAP utilizam
\textit{crawlers} que analisam estaticamente o HTML em busca de âncoras
\texttt{<a>}, percebendo uma SPA como uma única página opaca
\cite{brightsec2024}.

O problema é amplificado por uma segunda dimensão: a incapacidade de
interpretação semântica. Vulnerabilidades como Lógica de Negócio
(\textit{Business Logic Vulnerabilities} --- BLVs), IDOR
(\textit{Insecure Direct Object Reference}), XSS armazenado em campos
dinâmicos, adulteração de JWT e exposição de dados sensíveis em
\textit{endpoints} não documentados não decorrem de padrões sintáticos
detectáveis por expressões regulares ou por injeção de \textit{payloads}
clássicos \cite{portswigger2024, karadbhajne2025}. Elas exigem
compreensão do \textit{propósito} da aplicação: saber que uma quantidade
de produto deve ser positiva, que um carrinho pertence a um único usuário,
ou que uma rota de administração não deve ser acessível por clientes.

% --- Parte 3: Tese / Solução ----------------------------------
\textbf{Nossa tese é que LLMs agênticos, equipados com automação de
navegadores e orientados pelo padrão ReAct, são capazes de navegar
autonomamente os estados internos de uma SPA e detectar vulnerabilidades
que requerem raciocínio semântico sobre o domínio da aplicação},
superando estatisticamente as ferramentas DAST tradicionais nesse
subconjunto de ameaças. A abordagem proposta --- o \textbf{Agentic DAST}
--- não busca substituir \textit{scanners} tradicionais para
vulnerabilidades técnicas (SQLi, XSS por injeção direta), mas
complementá-los na detecção das vulnerabilidades semanticamente
complexas que escapam à automação sintática.

% --- Parte 4: Contribuições -----------------------------------
As contribuições deste trabalho são:

\begin{enumerate}\setlength{\itemsep}{2pt}
  \item \textbf{\textit{Framework} Agentic DAST} — arquitetura de quatro
  módulos (Percepção, Cognitivo, Executor, Auditoria) implementada em
  TypeScript, disponibilizada como \textit{software} de código aberto;

  \item \textbf{Biblioteca de \textit{Abuse Cases}} — oito casos de abuso
  predefinidos cobrindo as categorias BLV, \textit{Broken Access Control},
  XSS, JWT, \textit{Sensitive Data Exposure} e SQL Injection, modelados
  com o padrão Screenplay;

  \item \textbf{Protocolo de avaliação empírica} — experimento controlado
  comparativo com \textit{ground truth} derivado do catálogo oficial
  OWASP Juice Shop, usando as métricas \textit{Recall}, \textit{Precision}
  e F1-\textit{Score} por categoria de vulnerabilidade.
\end{enumerate}

O restante deste artigo está organizado da seguinte forma: a
Seção~\ref{sec:relacionados} discute os trabalhos relacionados;
a Seção~\ref{sec:framework} descreve a arquitetura do Agentic DAST;
a Seção~\ref{sec:avaliacao} apresenta o protocolo experimental;
a Seção~\ref{sec:conclusao} conclui o trabalho.

% =============================================================
\section{Trabalhos Relacionados}
\label{sec:relacionados}
% =============================================================

% --- DAST em SPAs -------------------------------------------
\subsection{DAST Tradicional e Limitações em SPAs}

O OWASP ZAP \cite{owaspzap} é a ferramenta DAST de referência no meio
acadêmico e na indústria. Sua arquitetura de \textit{spider} foi
projetada para o paradigma das \textit{Multi-Page Applications} (MPAs),
onde cada interação HTTP resulta em uma nova página HTML analisável
estaticamente. \citeonline{brightsec2024} demonstraram empiricamente que
o ZAP identifica menos de 20\% dos \textit{endpoints} de uma SPA Angular
em comparação com uma suíte de testes E2E que navega os estados
JavaScript. \citeonline{stackhawk2024} propõem integrar o ZAP a scripts
Playwright pré-gravados como \textit{baseline} --- abordagem que
instrumenta o \textit{crawler} com conhecimento humano prévio, mas não
oferece descoberta autônoma.

\citeonline{mdpiBLV2025} realizaram uma revisão sistemática sobre
ferramentas de detecção de BLVs, concluindo que nenhuma ferramenta
automatizada disponível detecta BLVs de forma confiável sem intervenção
humana na especificação das regras de negócio.

% --- LLMs em Segurança --------------------------------------
\subsection{LLMs Aplicados à Segurança Ofensiva}

\citeonline{yang2025hybrid} demonstraram que LLMs com
\textit{constraint solving} geram \textit{payloads} que cobrem caminhos
de execução profundos inacessíveis a \textit{fuzzers} por mutação
aleatória em protocolos de comunicação. \citeonline{lsast2024}
propuseram o LSAST, um sistema de análise estática de segurança guiado
por LLM com capacidade de identificar vulnerabilidades lógicas em código
PHP. Ambos os trabalhos operam, contudo, sobre código-fonte ou protocolos
estáticos --- não sobre aplicações web em execução.

\citeonline{pentest2024llm} investigaram o uso de LLMs como assistentes
de \textit{pentest} manual, demonstrando que modelos como GPT-4 conseguem
raciocinar sobre vulnerabilidades quando fornecidos com informações
estruturadas da aplicação alvo. O Agentic DAST avança esse trabalho ao
eliminar a intermediação humana: o LLM recebe o estado da aplicação em
tempo real e executa as ações diretamente.

% --- ReAct e Agentes Autônomos ------------------------------
\subsection{Padrão ReAct e Agentes de Software}

\citeonline{yao2023react} introduziram o padrão ReAct
(\textit{Reasoning + Acting}), no qual o LLM intercala passos de
raciocínio (\textit{thought}), ações (\textit{action}) e observações
do ambiente (\textit{observation}), produzindo trajetórias de decisão
interpretáveis e auditáveis. \citeonline{wei2022chain} demonstraram que
\textit{Chain-of-Thought prompting} habilita LLMs a decompor problemas
complexos em sub-etapas, elevando significativamente a acurácia em
tarefas de raciocínio de múltiplas etapas.

% --- Padrão Screenplay --------------------------------------
\subsection{Padrão Screenplay em Automação de Testes}

O \textit{Screenplay Pattern} \cite{serenityscreenplay}, evolução do
\textit{Page Object Model} (POM), modela testes em termos de Atores,
Habilidades, Tarefas e Interações atômicas --- desacoplando a intenção
(\textit{Task}) da implementação (\textit{Interaction}). Isso permite
que o agente LLM raciocine em termos semânticos (``Tentar login com
injeção SQL'') sem conhecer os seletores CSS específicos, que são
resolvidos em tempo de execução pelo Executor.

% --- Tabela Comparativa -------------------------------------
\subsection{Comparação com Trabalhos Próximos}

A Tabela~\ref{tab:relacionados} compara o Agentic DAST com os trabalhos
mais próximos em termos de objetivo, abordagem e cobertura.

\begin{table*}[!htb]
\centering
\caption{Comparação com trabalhos relacionados. ``Auto.''~=~detecção
autônoma (sem scripts pré-gravados); ``Sem.''~=~raciocínio semântico;
``SPA''~=~suporte a SPAs com JavaScript.\label{tab:relacionados}}
\small
\begin{tabularx}{\textwidth}{lXcccc}
\toprule
\textbf{Trabalho} & \textbf{Abordagem} & \textbf{Auto.} & \textbf{Sem.} & \textbf{SPA} & \textbf{BLV} \\
\midrule
OWASP ZAP \cite{owaspzap}
  & \textit{Spider} + scan ativo & Sim & Não & Não & Não \\
ZAP + E2E \cite{stackhawk2024}
  & ZAP guiado por script humano & Não & Não & Sim & Não \\
LSAST \cite{lsast2024}
  & Análise estática com LLM & Sim & Sim & Não & Parcial \\
LLM Pentest \cite{pentest2024llm}
  & LLM assistindo humano & Não & Sim & Sim & Sim \\
Yang et al. \cite{yang2025hybrid}
  & Fuzzing + LLM (protocolos) & Sim & Sim & Não & Não \\
\midrule
\textbf{Agentic DAST (este trabalho)}
  & Agente LLM autônomo + SPA & \textbf{Sim} & \textbf{Sim} & \textbf{Sim} & \textbf{Sim} \\
\bottomrule
\end{tabularx}
\end{table*}

A tabela evidencia que nenhum trabalho anterior combina simultaneamente
autonomia total, raciocínio semântico, suporte a SPAs com JavaScript e
detecção de BLVs. O Agentic DAST é, ao nosso conhecimento, o primeiro
\textit{framework} a abordar todas essas dimensões de forma integrada.

% =============================================================
\section{Framework Proposto}
\label{sec:framework}
% =============================================================

\subsection{Visão Geral da Arquitetura}

O Agentic DAST é estruturado em quatro módulos com responsabilidades
bem definidas, que operam em sequência conforme a
Figura~\ref{fig:arquitetura}. O fluxo de execução ocorre em quatro
fases: (1)~\textit{Percepção}: o DOMExtractor navega a SPA e constrói
o \textit{AppMap}; (2)~\textit{Inicialização}: o ActionDispatcher
instancia o Actor com as \textit{Abilities}; (3)~\textit{Loop ReAct}:
o ReActAgent intercala Thought/Action/Observation com o LLM;
(4)~\textit{Relatório}: os \textit{Findings} são persistidos em JSON.

\begin{figure}[!htb]
\centering
\small
\begin{verbatim}
  ┌─────────────────────────────────────┐
  │          FASE 1: PERCEPÇÃO          │
  │  DOMExtractor (Puppeteer + CDP)     │
  │  → AppMap {rotas, endpoints, JWT}   │
  └──────────────┬──────────────────────┘
                 │ AppMap
  ┌──────────────▼──────────────────────┐
  │       FASE 3: LOOP REACT            │
  │  ReActAgent (Claude API)            │
  │  Thought → Action → Observation     │
  │  ↕ ActionDispatcher                 │
  │  ↕ BrowseTheWeb (Puppeteer)         │
  │  ↕ HttpClient (Axios)               │
  │  ↕ OWASP ZAP (proxy passivo)        │
  └──────────────┬──────────────────────┘
                 │ Findings
  ┌──────────────▼──────────────────────┐
  │       FASE 4: RELATÓRIO             │
  │  report-<timestamp>.json            │
  └─────────────────────────────────────┘
\end{verbatim}
\caption{Arquitetura do Agentic DAST. O loop ReAct opera sobre o
AppMap gerado na fase de Percepção; o ZAP audita passivamente todo
o tráfego gerado pelo Puppeteer.}
\label{fig:arquitetura}
\end{figure}

\subsection{Módulo de Percepção}

O \texttt{DOMExtractor} é implementado com Puppeteer, que controla o
Chromium via \textit{Chrome DevTools Protocol} (CDP). A escolha do CDP
sobre o Selenium WebDriver é justificada pela necessidade de
intercepção de rede de baixo nível: o CDP permite pausar, inspecionar
e modificar requisições HTTP antes de chegarem ao renderizador, acesso
impossível com o Selenium \cite{lightpanda2024}.

O módulo realiza: (i)~navegação por todas as rotas Angular registradas
no \textit{router}; (ii)~extração de formulários, botões e
metadados ARIA de cada estado; (iii)~interceptação de chamadas
\texttt{/api/} e \texttt{/rest/} via CDP; (iv)~login automatizado
e extração do JWT do \texttt{localStorage}. O resultado é o
\textbf{AppMap} --- um objeto JSON estruturado que é a entrada para
o Módulo Cognitivo.

\subsection{Módulo Cognitivo}

O \texttt{ReActAgent} implementa o \textit{loop} ReAct sobre a API da
Anthropic (modelo Claude). O \texttt{SYSTEM\_PROMPT}, marcado com
\textit{Prompt Caching} da Anthropic, instrui o LLM sobre: (i)~seu
papel como pesquisador de segurança autorizado; (ii)~as seis categorias
de vulnerabilidades a investigar; (iii)~as oito ações disponíveis e
seus formatos JSON. O caching reduz o custo de \textit{tokens} em até
90\% nas iterações subsequentes \cite{anthropicpromptcaching}.

O histórico de conversa cresce iterativamente: AppMap inicial,
seguido de ciclos (\textit{thought} + \textit{action}) /
\textit{observation}. O LLM decide autonomamente qual ação executar
em cada etapa com base nas observações acumuladas.

\subsection{Módulo Executor}

O Executor é implementado com o \textit{Screenplay Pattern}
\cite{serenityscreenplay} em três camadas:

\begin{itemize}\setlength{\itemsep}{2pt}
  \item \textbf{Actor} --- o agente de segurança, que agrega as
  habilidades via interface fluente
  (\texttt{Actor.named().withBrowser().withHttpClient()});

  \item \textbf{Abilities} --- \texttt{BrowseTheWeb} (Puppeteer +
  CDP) e \texttt{HttpClient} (Axios com \texttt{validateStatus:
  () => true}, que nunca lança exceção em 4xx/5xx);

  \item \textbf{ActionDispatcher} --- mapeia cada
  \texttt{ParsedAction} (tipo discriminado TypeScript) para a
  \textit{Ability} correspondente, com propagação automática do JWT
  extraído para o \texttt{HttpClient}.
\end{itemize}

As oito ações disponíveis ao LLM são: \texttt{navigate},
\texttt{fill\_field}, \texttt{click\_element}, \texttt{read\_page},
\texttt{make\_request}, \texttt{extract\_tokens},
\texttt{report\_finding} e \texttt{finish}.

Complementando o \textit{loop} ReAct, a \textbf{Biblioteca de Abuse
Cases} executa oito testes estruturados após o \textit{loop},
garantindo cobertura mínima das vulnerabilidades conhecidas do Juice
Shop. A Tabela~\ref{tab:abusecases} lista os \textit{Abuse Cases}
implementados.

\begin{table}[!htb]
\centering
\caption{Biblioteca de \textit{Abuse Cases} implementados, com
categorias OWASP e severidade esperada.\label{tab:abusecases}}
\small
\begin{tabularx}{\columnwidth}{lXl}
\toprule
\textbf{\textit{Abuse Case}} & \textbf{Categoria} & \textbf{Sev.} \\
\midrule
NegativeQty     & BLV — Lógica de Negócio          & HIGH     \\
ZeroStarReview  & BLV — Validação de Entrada        & MEDIUM   \\
IDOR            & Broken Access Control (A01)       & HIGH     \\
AdminRoute      & Broken Access Control (A01)       & CRITICAL \\
XSS             & Injection — XSS (A03)             & HIGH     \\
JWT             & Falhas Criptográficas (A02)       & CRITICAL \\
SensitiveData   & Misconfiguration (A05)            & HIGH     \\
SQLInjection    & Injection — SQLi (A03)            & CRITICAL \\
\bottomrule
\end{tabularx}
\end{table}

\subsection{Módulo de Auditoria}

O OWASP ZAP opera em modo \textit{daemon} (via Docker) como proxy
reverso na porta 8888. Todo o tráfego gerado pelo Puppeteer durante a
fase de Percepção é roteado pelo ZAP, que registra passivamente a
superfície de ataque. Após o \textit{loop} ReAct, o
\texttt{ZapClient} consulta a API REST do ZAP para recuperar os
alertas acumulados, filtrando aqueles com confiança zero (falsos
positivos prováveis). Os alertas são convertidos para o formato
\texttt{Finding} do \textit{framework} e consolidados no relatório
final junto com os \textit{Findings} do agente e dos \textit{Abuse Cases}.

% =============================================================
\section{Avaliação}
\label{sec:avaliacao}
% =============================================================

Esta seção descreve o protocolo experimental adotado para validar o
Agentic DAST. Os resultados quantitativos serão apresentados na versão
final deste artigo (TCC~II), após a execução dos experimentos previstos
para o segundo semestre de 2026.

\subsection{Objetivos e Perguntas de Pesquisa}

O experimento é estruturado em torno de três Perguntas de Pesquisa (RQs):

\begin{description}\setlength{\itemsep}{2pt}
  \item[\textbf{RQ1}] \textit{O Agentic DAST detecta mais
  vulnerabilidades que o OWASP ZAP com \textit{spider} tradicional?}
  Esta RQ avalia o Recall geral e por categoria de vulnerabilidade.

  \item[\textbf{RQ2}] \textit{O agente LLM descobre mais
  \textit{endpoints} e estados de UI que scanners tradicionais?}
  Esta RQ avalia a capacidade de navegação semântica do framework.

  \item[\textbf{RQ3}] \textit{Os \textit{Abuse Cases} predefinidos
  alcançam maior taxa de confirmação de vulnerabilidades que abordagens
  de \textit{fuzzing} aleatório sobre os mesmos campos?}
  Esta RQ avalia a qualidade semântica dos testes estruturados.
\end{description}

\subsection{\textit{Setup} Experimental}

O ambiente de testes é composto por três contêineres Docker
orquestrados via Docker Compose:

\begin{itemize}\setlength{\itemsep}{2pt}
  \item \textbf{Juice Shop} (\texttt{bkimminich/juice-shop:latest}) ---
  a SPA Angular alvo, exposta na porta 3000;
  \item \textbf{OWASP ZAP} (\texttt{ghcr.io/zaproxy/zaproxy:stable})
  --- \textit{daemon} na porta 8080 (API) e 8888 (proxy);
  \item \textbf{Agente} --- processo Node.js 20 com TypeScript,
  consumindo a API Claude (\texttt{claude-sonnet-4-6}).
\end{itemize}

Cada cenário é executado $n = 3$ rodadas independentes; o banco de
dados do Juice Shop é reinicializado entre rodadas para garantir estado
inicial idêntico. A temperatura do LLM é fixada em $T = 0{,}3$ para
reduzir a variância entre rodadas.

Os três cenários comparativos são:

\begin{itemize}\setlength{\itemsep}{2pt}
  \item \textbf{Cenário~A} --- OWASP ZAP com \textit{spider} padrão,
  sem qualquer instrumentação adicional (\textit{baseline} tradicional);
  \item \textbf{Cenário~B} --- OWASP ZAP guiado por suíte de testes E2E
  do ``caminho feliz'' pré-gravada (\textit{baseline} com navegação
  humana);
  \item \textbf{Cenário~C} --- Agentic DAST em modo \textit{black-box}
  autônomo, sem conhecimento prévio da estrutura da aplicação
  (cenário experimental).
\end{itemize}

\subsection{\textit{Ground Truth} e Métricas}

A Matriz de \textit{Ground Truth} (GT) é construída a partir do catálogo
oficial de desafios do OWASP Juice Shop \cite{owaspjuiceshop},
selecionando os challenges das categorias investigadas (BLV, Broken
Access Control, XSS, JWT, Sensitive Data, Injection). Cada challenge
é classificado manualmente antes da execução dos experimentos, garantindo
gabarito independente dos resultados.

As métricas coletadas para cada cenário são:

\begin{table}[!htb]
\centering
\caption{Métricas do experimento. VP~=~Verdadeiro Positivo,
FP~=~Falso Positivo, FN~=~Falso Negativo.\label{tab:metricas}}
\small
\begin{tabularx}{\columnwidth}{lX}
\toprule
\textbf{Métrica} & \textbf{Definição} \\
\midrule
\textit{Recall}    & $\text{VP} / (\text{VP} + \text{FN})$ \\
\textit{Precision} & $\text{VP} / (\text{VP} + \text{FP})$ \\
F1-\textit{Score}  & $2 \cdot (\text{Prec} \cdot \text{Rec}) /
                     (\text{Prec} + \text{Rec})$ \\
Cobertura          & N.º de \textit{endpoints} e estados UI descobertos \\
\bottomrule
\end{tabularx}
\end{table}

\subsection{Hipóteses Testáveis}

\begin{description}\setlength{\itemsep}{2pt}
  \item[\textbf{H1}] $\text{Recall}(C) > \text{Recall}(A)$ com
  significância estatística (teste de Wilcoxon, $p < 0{,}05$).
  \item[\textbf{H2}] $\text{Endpoints}(C) > \text{Endpoints}(A)$ e
  $\text{Endpoints}(C) > \text{Endpoints}(B)$.
  \item[\textbf{H3}] A taxa de confirmação dos \textit{Abuse Cases}
  é superior à do \textit{fuzzing} aleatório sobre os mesmos campos.
\end{description}

A análise estatística utiliza o teste de Wilcoxon para amostras
emparelhadas com $\alpha = 0{,}05$, adequado ao tamanho de amostra
reduzido ($n = 3$) sem pressuposição de normalidade
\cite{kitchenham2002}.

\subsection{Ameaças à Validade}

\textbf{Validade interna:} a natureza probabilística dos LLMs pode
causar variação entre rodadas. Controlamos com temperatura fixa
($T = 0{,}3$), mesma versão do modelo e reinicialização do banco de
dados. \textbf{Validade externa:} os resultados são obtidos sobre o
Juice Shop, uma aplicação sintética; a generalização para sistemas reais
de produção requer estudos adicionais em ambientes de \textit{staging}.
\textbf{Validade de construção:} a Matriz GT é derivada do catálogo
oficial OWASP, minimizando subjetividade; vulnerabilidades ambíguas são
excluídas do conjunto de avaliação.

% =============================================================
\section{Conclusão}
\label{sec:conclusao}
% =============================================================

Este artigo apresentou o \textbf{Agentic DAST}, um \textit{framework}
que endereça a lacuna crítica entre as capacidades das ferramentas DAST
tradicionais e as demandas de segurança de SPAs modernas. O problema
é duplo: as SPAs tornam a superfície de ataque inacessível a
\textit{crawlers} estáticos, e as vulnerabilidades semanticamente
complexas (BLV, IDOR, JWT) são invisíveis a ferramentas baseadas em
padrões sintáticos.

Nossa tese --- de que LLMs agênticos com automação de navegadores podem
preencher essa lacuna --- é materializada em uma arquitetura de quatro
módulos integrados: Percepção via CDP, raciocínio ReAct sobre a API
Claude, execução via Screenplay e auditoria passiva via OWASP ZAP.
A biblioteca de oito \textit{Abuse Cases} predefinidos garante cobertura
mínima e reprodutível das categorias de vulnerabilidade investigadas.

\textbf{Limitações:} o \textit{framework} atual é validado exclusivamente
sobre o OWASP Juice Shop; a latência e o custo das chamadas à API do LLM
tornam o Agentic DAST complementar, e não substituto, ao ZAP em
\textit{pipelines} de CI/CD de alta frequência. O número reduzido de
rodadas ($n = 3$) limita o poder estatístico das comparações.

\textbf{Trabalhos futuros:} (i)~ampliar a validação para SPAs de
produção em ambientes de \textit{staging} corporativo; (ii)~investigar
modelos LLM locais (Llama, Mistral) como alternativas de menor custo e
maior privacidade; (iii)~estender o \textit{Screenplay Pattern} com
suporte a múltiplos atores simultâneos (vítima + atacante) para testes
de vulnerabilidades de controle de acesso mais sofisticados;
(iv)~integrar o \textit{framework} como \textit{plugin} nativo do
OWASP ZAP.

Os resultados empíricos do experimento comparativo (Cenários A, B e C)
serão reportados na continuação deste trabalho (TCC~II), com potencial
de impacto direto em \textit{pipelines} de DevSecOps na indústria
brasileira e na literatura científica sobre DAST agêntico em SPAs.

% =============================================================
% Referências
% =============================================================

\bibliography{referencias}

\end{document}
